\chapter{Analyse et conception}

\section{Étude de l'existant}
\TODO{Décrire les outils/processus existants avant SRS pour le cadrage de projets
similaires au sein de la Fondation OCP, si une étude a été menée. Sinon, indiquer
explicitement qu'aucune étude formelle de l'existant n'a été réalisée dans le cadre de
ce stage.}

\section{Analyse des besoins}
Les besoins ont été structurés autour d'un modèle de spécification centralisé (le
\emph{Knowledge Graph}), qui organise les réponses de l'entretien par domaine~:
parties prenantes, utilisateurs, coopératives, bénéficiaires, exigences fonctionnelles,
règles métier, workflows, documents, KPI, ODD, ESG, sécurité, données, intégrations,
architecture, infrastructure, notifications, reporting, tests et déploiement.

\section{Acteurs et rôles}
\begin{itemize}
    \item \textbf{Administrateur / consultant SRS} — mène l'entretien de cadrage.
    \item \textbf{Partie prenante interrogée} — représentant métier de la Fondation OCP
    ou de l'Axe Économie Sociale et Solidaire.
    \item \textbf{Équipe de développement (destinataire des livrables)} — reprend le
    Cahier des Charges et la conception pour développer la future plateforme.
\end{itemize}

\section{Fonctionnalités principales de SRS}
\begin{itemize}
    \item Entretien structuré (Classic Questionnaire hors ligne ou IA adaptative).
    \item Moteur de Knowledge Graph avec suivi de complétude et de dépendances entre
    concepts.
    \item Moteur de génération d'exigences (fonctionnelles, non fonctionnelles,
    sécurité, techniques).
    \item Moteur d'analyse de sécurité (score de sécurité, checklist, matrice RBAC,
    registre des risques, analyse d'impact vie privée).
    \item Moteur de revue (règles de complétude avant export final).
    \item Générateur de documents multi-formats (JSON, Markdown, HTML, DOCX, PDF).
\end{itemize}

\section{Architecture générale}
SRS suit une architecture classique en trois couches~:

\begin{itemize}
    \item \textbf{Frontend} — application Next.js (React) assurant l'assistant
    d'entretien, les tableaux de bord et la consultation des livrables générés.
    \item \textbf{Backend} — API FastAPI (Python) exposant les routes projet, entretien,
    knowledge graph, exigences, sécurité, revue et documents.
    \item \textbf{Base de données} — PostgreSQL en production (SQLite en environnement
    de test), via SQLAlchemy asynchrone.
\end{itemize}

Il n'existe volontairement \textbf{aucune couche d'authentification} dans SRS~: chaque
projet appartient à un contexte interne unique (« Fondation OCP »), l'outil étant destiné
à un usage interne restreint pendant la phase de cadrage plutôt qu'à un accès multi-
organisation.

\section{Conception de la solution}
La conception distingue explicitement, à chaque étape de la génération documentaire~:
les informations \emph{confirmées} par une réponse d'entretien, les
\emph{hypothèses} formulées par le système, et les points \emph{à confirmer} avec les
parties prenantes concernées (par exemple les exigences légales/RGPD ou Loi~09-08, qui
sont systématiquement étiquetées comme à valider avec le responsable juridique/DPO).

\section{Diagramme de conception du MVP}
Conformément au principe « un seul diagramme de conception MVP compréhensible par des
publics techniques et non techniques », SRS génère dynamiquement un diagramme unique
(format Mermaid) représentant les acteurs, le frontend, le backend, les modules
effectivement justifiés par les exigences confirmées, la base de données, et les
contrôles de sécurité applicables — sans ajout systématique de briques non justifiées
(microservices, Kubernetes, files de messages, etc.).

\begin{figure}[h]
    \centering
    \TODO{Insérer une capture d'écran ou un rendu du diagramme MVP généré pour un
    projet réel, une fois disponible.}
    \caption{Exemple de diagramme de conception MVP généré par SRS}
\end{figure}

\section{Flux principaux}
\subsection{Flux d'entretien}
Démarrage de session → question adaptative ou issue du catalogue de concepts →
interprétation de la réponse → mise à jour du Knowledge Graph → génération incrémentale
des exigences associées → question suivante ou fin d'entretien.

\subsection{Flux de génération documentaire}
Sur demande, le générateur de documents assemble le contenu à partir de l'état courant
du Knowledge Graph, des exigences et — si disponible — de l'analyse de sécurité, puis
l'exporte dans le format demandé (JSON, Markdown, HTML, DOCX, PDF).

\subsection{Flux de revue}
Avant export final, le moteur de revue vérifie la complétude (parties prenantes,
exigences, sécurité, architecture, tests) et distingue les éléments confirmés, supposés
et restant à confirmer.

\section{Choix technologiques}
\begin{longtable}{|l|l|}
\hline
\textbf{Composant} & \textbf{Technologie} \\
\hline
Backend & FastAPI, SQLAlchemy (async), Alembic \\
Base de données & PostgreSQL (prod), SQLite (tests) \\
Frontend & Next.js, React, Tailwind CSS \\
IA (entretien adaptatif) & Google Gemini, via une abstraction de client IA \\
Export documentaire & python-docx, ReportLab, markdown2 \\
Tests backend & pytest, pytest-asyncio \\
\hline
\end{longtable}
